\section{Conditional finite-data reconstruction}\label{sec:stability}
\subsection{The prior class and the loss}
Fix a compact class of physical presentations with at most $J$ records
and at most $2J$ obstacle representatives. All selected channels have
uniform positive separation, isolation and adjacent free-flight margins;
gaps, curvature, free area and period covolume have positive lower bounds.
Diameters, gaps, areas and some lattice bases and their inverses have
fixed upper bounds. Each selected onset belongs to a known bounded bracket,
and queries in a fixed collar beyond its upper endpoint are allowed.
These are local itinerary brackets, not exact onsets or hidden contacts.

Steiner-centered support functions extend holomorphically to
$|\operatorname{Im}z|<\rho$ with a common bound $M$. For constants
$\eta,\nu>0$ require
\begin{align}
 \|p_i(\cdot+\phi)-p_i\|_2
   &\ge\eta\dist(\phi,2\pi\Z),                 \label{eq:rot-margin}\\
 \|p_i(-\cdot+\phi)-p_i\|_2&\ge\eta,         \label{eq:ref-margin}\\
 \inf_{O\in O(2)}\|p_i-O\cdot p_j\|_2&\ge\nu\quad(i\ne j).
                                                        \label{eq:shape-margin}
\end{align}
The norms are on the real normal-angle circle. Every fixed finite
asymmetric, shape-separated family admits such margins. Near zero rotation,
the first follows from $\|p_i'\|_2>0$; elsewhere and for reflections or
distinct shapes, compactness of $O(2)$ and absence of equality give the
positive bounds. The bounds persist in a small compact neighborhood.
They are quantitative priors, not a supplied shape catalogue.

Uniform analytic inverse-function and stationary-branch constructions,
with positive curvature and twist, give common local contact graphs,
active endpoint rectangles and fixed derivative bounds for the densities.
Indeed the stationary derivative is bounded away from zero at the normal
orbit; compactness and the uniform complex graph bounds retain a common
smaller neighborhood. Thus one may impose any fixed
$C^{2,\beta}$ bound there, $0<\beta\le1$. The physical $C^3$ and all
other hypotheses of Proposition~\ref{prop:local-stable} are included.
The family has the connected coverage and primitive-cycle property of
Theorem~\ref{thm:stat-main}; the graph and cycle pair are not inputs.

Here is the precise loss. In each comparison allow a permutation of
representatives, one common Euclidean isometry, and changes of period
representatives along a spanning tree of the recovered incidence graph.
Allow also unmarked lattice bases whose norms and inverse norms are at
most a fixed $B$. Take the infimum of
\begin{equation}\label{eq:loss}
       \max_i\|p_{C_i}-p_{\widetilde C_i}\|_{C^2(S^1)}
                            +\|L-\widetilde L\|
\end{equation}
over these bounded presentations, applying the same orthogonal map to
the lattice and all bodies. Representatives are in a bounded root-tree
gauge. Choose $B$ large enough in terms of the prior to include every
observed primitive cycle basis: its columns are bounded finite sums of
channel placements and its determinant is bounded below by $V$.
We call this the bounded-presentation loss $\mathcal E$. It asserts
closeness of actual boundary images and an unmarked period lattice, not
independent per-edge placements. The proof establishes the explicit
comparison \eqref{eq:loss}; no metric property of this infimum is needed.
The compact physical class can equivalently be specified using these
bounded presentations modulo their finite relabellings and gauges.

\subsection{From cell probabilities to density derivatives}
Choose centered rectangles $Q_0\Subset Q$ and square grids invariant
under simultaneous reversal. A histogram records the cells of $Q$,
success outside $Q$, and failure. A single phase preparation gives one
multinomial outcome, not a separate trial for every cell.

\begin{lemma}[Cell reconstruction]\label{lem:bins}
For densities uniformly bounded in $C^{2,\beta}(Q)$, cell probabilities
on a grid of spacing $h$ determine a linear reconstruction with
\begin{equation}\label{eq:bins}
     \|\widehat f-f\|_{C^2(Q_0)}
           \le C(h^\beta+\epsilon h^{-4})
\end{equation}
when each probability has error at most $\epsilon$. The same bound
holds for the distance between two physical densities whose cell
probabilities differ by at most $\epsilon$. There are $O(h^{-2})$ cells.
\end{lemma}
\begin{proof}
The averages of $1,z,z^2$ on unit cells centered at $i=-1,0,1$ form
the invertible matrix with rows $(1,i,i^2+1/12)$. Its tensor square
recovers a coordinatewise quadratic polynomial from nine cell averages.
For a $C^{2,\beta}$ function the total-degree two Taylor polynomial
has remainder $O(h^{2+\beta})$, and it is reproduced exactly by this
operator. Its derivative errors of orders $j\le2$ are
$O(h^{2+\beta-j})$ on a fixed enlarged block. A probability error
$\epsilon$ gives average-density error $\epsilon h^{-2}$ and derivative
error $O(\epsilon h^{-2-j})$. Patch the polynomials with a smooth
partition of unity whose $j$th derivatives are $O(h^{-j})$ and whose
overlap number is bounded. Leibniz's rule gives the same bounds. The
larger rectangle supplies all boundary stencils for $Q_0$.
Apply this construction to each of two densities and subtract to obtain
the comparison assertion.
\end{proof}

We will use the smaller variance of a cell indicator. If $N$ independent
preparations are made at an active time, $p_B\le Ch^2$ and Bernstein's
inequality gives, simultaneously for all cells with probability at least
$1-\delta$, the bound
\begin{equation}\label{eq:bernstein}
 |\widehat p_B-p_B|\le C(h\sqrt{r_N}+r_N),\qquad
 r_N=\frac{\log(CJN/\delta)}{N},
\end{equation}
provided the grid has at most $CN$ cells. Dependence among the indicators
of one multinomial sample does not affect a union bound. Together with
Lemma~\ref{lem:bins}, this gives
\begin{equation}\label{eq:local-noise}
           e_N=C\{h^\beta+\sqrt{r_N}h^{-3}+r_Nh^{-4}\}.
\end{equation}
This is a variance-aware refinement of the absolute-error histogram bound,
not a lower-bound or optimality assertion.

\subsection{Analytic continuation and shape synchronization}
\begin{lemma}[A uniform continuation estimate]\label{lem:continuation}
Let $p,\widetilde p$ be periodic holomorphic functions bounded on a
common strip of positive width. If they differ by at most $e$ on a
fixed positive-length real interval, then
\[
             \|p-\widetilde p\|_{C^2(S^1)}\le Ce^\theta
\]
for some $C>0$ and $0<\theta<1$ depending on the strip, bound and
interval length. The constants are uniform under translation of the
interval on the circle.
\end{lemma}
\begin{proof}
Slit a narrower periodic strip cylinder along a closed shorter interval.
The harmonic measure $\omega$ of the slit is positive in this connected
domain, equals one on the slit, and zero on the two outer boundary
circles. Its minimum on a further narrower closed cylinder, extended
by one at the slit, is some $\theta>0$. The slit endpoints are regular
Dirichlet boundary points. The subharmonic maximum principle gives
$|p-\widetilde p|\le(2M)^{1-\omega}e^\omega$. Zeros are handled by
regularization. Cauchy estimates on fixed disks in the narrower cylinder
give the two derivative bounds on the real circle. For $e=0$ use the
identity theorem. This is the classical two-constants mechanism; its
geometric loss is consistent with the analysis in \cite{Trefethen}.
\end{proof}

\begin{lemma}[Stable incidence and geometric cycles]\label{lem:sync-stable}
Suppose two physical collections in the prior class can be paired record
by record, with an independent coherent reversal on each record, so their
active densities have $C^2$ discrepancy at most $e$. For sufficiently
small $e$ they have isomorphic incidence multigraphs and the same number
of obstacle classes. In a common root-tree gauge their corresponding
representative images, cycle displacements and calibrated covolumes
have error at most $Ce^\theta$.
\end{lemma}
\begin{proof}
Proposition~\ref{prop:local-stable} gives $Ce$ error for each pair of
actual arcs in its contact frame. Positive curvature converts them to
support functions on a common normal-angle interval. The normal-angle
map has derivative equal to curvature, bounded below; its inverse and
the support function through two derivatives depend Lipschitzly on the
arc, using the physical $C^3$ prior. Lemma~\ref{lem:continuation} gives
$\Delta=Ce^\theta$ error for each complete pair image.

Steiner centering is a bounded linear operation on support functions.
If two occurrences represent the same shape in one table, their partners
in the other have congruence distance at most $C\Delta$. They must belong
to the same class by \eqref{eq:shape-margin}. Conversely two distinct
classes cannot become congruent when $C\Delta<\nu/4$. Thus the equality
pattern of all occurrences, and hence the incidence multigraph, agrees.

Two approximate matchings of one source image to a placed representative
differ by an orthogonal transformation moving its centered support
function by $O(\Delta)$ in $L^2$. The reflection margin excludes the
incorrect component; the rotation margin bounds its angle by
$C\Delta/\eta$. Uniform third derivatives of support functions control
the induced $C^2$ error, and Steiner points control translations.
Induction along a fixed spanning tree aligns all representative images
with error $C\Delta$. The same argument for a non-tree pair and
\eqref{eq:cycle-displacement} give $C\Delta$ error in each $d_e$.
Finally the area error in the local inverse is $Ce$, and
\eqref{eq:support-area} bounds each obstacle-area error by $C\Delta$.
Summing each incidence class once proves the covolume estimate.
\end{proof}

\begin{lemma}[Integer locking for physical candidates]\label{lem:integer-lock}
Let two physical candidates have corresponding independent cycle matrices
$D,\widetilde D$ and covolumes $V,\widetilde V$ in the uniform bounded
class. If these quantities are sufficiently close, their calibrated
indices are the same integer $n$. If $n=1$, $D$ and $\widetilde D$
are bases of the two period lattices. More generally, corresponding
additional cycle vectors satisfy
\[
          nD^{-1}d_e=n\widetilde D^{-1}\widetilde d_e\in\Z^2
\]
when their discrepancy is sufficiently small, with the threshold also
depending on the bounded index.
\end{lemma}
\begin{proof}
The determinant-to-covolume ratio is Lipschitz on the bounded set with
$V$ bounded below. For physical candidates both ratios are positive
integers. A difference less than $1/2$ forces equality. Index one means
that the cycle subgroup already is the period lattice. At general index
$n$, the finite quotient $\Lambda/D\Z^2$ has order $n$, so
$n\Lambda\subset D\Z^2$. Thus the displayed vectors are integer.
Their difference is small by the inverse-matrix identity and the stated
bounds. An $\ell^\infty$ difference less than $1/2$ forces equality.
\end{proof}

The lemma applies to exact physical candidates compatible with uncertain
observations. It is not an instruction to take the integer span of noisy
real vectors, which need not be discrete, or to round uncalibrated real
numbers without an error bound. This distinction is important when no
integer copy labels are supplied.

\subsection{The pilot, estimator and preparation budget}
\begin{proof}[Proof of Theorem~\ref{thm:stat-main}]
Uniform local geometry gives $d_*,c,C>0$ such that
$cd^2\le\mu_{T_0+d}(\R^2)\le Cd^2$ for $0<d<d_*$, and zero mass
before the selected onset. One way to see the uniform quadratic order
is to scale the positive quadratic minimum of $W-T_0$ by $\sqrt d$
in \eqref{eq:density}; the two endpoint dimensions contribute $d$ and
the residual factor another $d$. Smooth Morse coordinates and the
positive twist make the two-sided bounds uniform.

Choose once and for all $\lambda>0$ with
$\sqrt{4\lambda/c}<d_*/32$, and a grid spacing less than $d_*/32$.
Query every onset bracket on this grid, including a short collar beyond
its endpoint, and let $\tau$ be the first reported mass exceeding
$2\lambda$. If the pilot errors are less than $\lambda/4$, no pre-onset
point can trigger; a grid point at offset between
$\sqrt{4\lambda/c}$ and that number plus one grid step must trigger.
Hence $0<\tau-T_0<d_*/16$. This uses no global monotonicity of the
three-impact probability. Use
$T_1=\tau+d_*/4$ and $T_2=\tau+d_*/2$. A uniform small rectangle
with $W-T_0\le d_*/8$ is strictly active at both times. All local
margins are positive. The pilot has a fixed prior-dependent number of
times per record, not a growing number as $N$ increases.

At each pilot time use $N$ fresh preparations, and at the two chosen
histogram times use fresh groups of $N$ preparations. Conditional on the
pilot, the histogram times are fixed and \eqref{eq:bernstein} applies.
Bernoulli concentration at the pilot and a union bound over all stages
and cells give pilot error $C\sqrt{r_N}$ and cell error
$C(h\sqrt{r_N}+r_N)$ with total failure probability at most $\delta$.
All preparations, including no-event outcomes, are charged.

For precision define an estimator over the compact physical prior, not
over arbitrary fitted density functions. Minimize the maximum discrepancy
from the retained pilot and histogram probabilities, divided by their
respective error tolerances, allowing a permutation of records and one
common reversal per record. Take an approximate minimizer from a fixed
countable dense family, within one of the infimum, using a first-eligible
index rule. The rule is Borel. The probability functions are continuous:
outside phase-null grazing, cell-boundary, patch-boundary and terminal-time
sets, collision records and intrinsic marks vary continuously. Uniform
separation bounds the number of collisions in a bounded queried interval,
and a fixed cell parametrization with $A$ bounded below permits dominated
convergence. Uniformly isolated contact branches also vary continuously.
Finite permutations and signs cause no obstruction. This justifies the
dense-family construction; no efficient search over analytic tables is
claimed.

On the concentration event the true table is admissible with normalized
discrepancy at most one. Thus the selected physical candidate has errors
bounded by fixed multiples of the stated tolerances. Keeping the pilot
in the discrepancy is essential: its first crossing also locates the
candidate's onset within the same fixed collar. For large $N$, candidate
pilot errors are still below $\lambda/4$, so both histograms are active
for it as well as for the truth. Lemma~\ref{lem:bins} gives the physical
pair discrepancy \eqref{eq:local-noise} with a changed constant.

Lemma~\ref{lem:sync-stable} recovers the incidence isomorphism and gives
$Ce_N^\theta$ errors in representatives, cycles and $V$. Choose, for the
comparison proof, a true spanning tree and cycle pair with calibrated
index one. The corresponding pair in the candidate has index one by
Lemma~\ref{lem:integer-lock}. Hence its cycle matrix is a period basis,
and \eqref{eq:loss} is at most $Ce_N^\theta$. The estimator was not
supplied this choice; the argument applies to every compatible candidate.

Finally take a centered square grid of spacing comparable to
$h=r_N^{1/(2\beta+6)}$. Then it has at most $CN$ cells and
\[
 h^\beta= r_N^{\beta/(2\beta+6)},\quad
 \sqrt{r_N}h^{-3}= r_N^{\beta/(2\beta+6)},\quad
 r_Nh^{-4}= r_N^{(2\beta+2)/(2\beta+6)}.
\]
The last term is smaller for $r_N<1$. There are at most $CJ$ queried
times, so the total is at most $CJN$ preparations. This proves
\eqref{eq:rate-main} and all discrete-stability assertions.
\end{proof}

The exponent is conditional on the analytic continuation geometry.
Near congruent shape classes or symmetric bodies, the incidence or
orientation thresholds deteriorate; without a uniform analytic strip,
exact image determination alone supplies no uniform noisy estimate.
At nonprimitive index the exact finite list remains valid, but a unique
estimate cannot be asserted merely by fitting arbitrarily fine histograms,
as Proposition~\ref{prop:ambiguity} demonstrates.
